\chapter{Introduction}
\label{chap:intro}

The development of this work is rooted in early exposure to embedded systems and aerial robotics during undergraduate studies at the Indian Institute of Technology (IIT) Jammu. Initial hands-on experience involved assembling and operating multiple drone platforms. While these systems enabled rapid prototyping, they largely relied on integrating existing components rather than building or understanding the system from first principles. This distinction highlighted a key limitation: the inability to fully control or adapt the underlying software stack.

During the later stages of the second year, an attempt was made to implement encrypted telemetry over low-frequency RF communication. Standard protocols such as MAVLink do not natively support secure communication, which required introducing additional hardware and software layers. The implemented solution involved routing telemetry data through UART to an ESP32 for AES-based encryption, transmitting over LoRa, and performing decryption on the receiving end before forwarding data to a ground station.

Although functional, this approach introduced significant complexity, latency, and development overhead. A feature that ideally should be integrated within the system required multiple intermediate components, increasing cost and points of failure. This experience highlighted a broader issue: existing drone software stacks are not designed for deep customization, lightweight integration of advanced features, or complete system control.

This realization led to a key insight. If implementing a relatively contained feature such as secure telemetry required substantial effort, the challenges faced by small teams or early-stage drone companies would be significantly greater. In practice, many such teams depend on large, externally developed platforms such as PX4 and ArduPilot, often modifying them to suit specific requirements. While these platforms provide flexibility, they introduce dependency on foreign-developed systems, increase system complexity, and require sustained engineering effort for customization and maintenance.

In addition, the lack of a widely adopted, publicly accessible indigenous drone flight stack in India presents both a challenge and an opportunity. Existing solutions are either highly abstracted or tightly coupled to specific hardware and software ecosystems, limiting their adaptability for emerging applications in robotics and autonomous systems. This dependency becomes particularly critical in domains such as defense, where reliance on external technology stacks is undesirable.

Another key limitation of current systems is hardware lock-in. Developing new hardware alongside a compatible firmware stack is a complex and resource-intensive process, especially for smaller teams. This barrier can slow innovation and restrict design flexibility.

Motivated by these challenges, the Vayu project was initiated in June 2025 with the goal of designing a complete flight control stack from first principles. The objective is to achieve vertical integration across all layers of the system, from hardware abstraction to high-level control logic, while maintaining modularity and performance.

A key design philosophy of this work is to minimize dependency on specific microcontroller families, external libraries, and monolithic frameworks. By maintaining control over all layers of the stack, the system aims to provide improved flexibility, portability, and long-term scalability.

The Vayu flight control stack is designed not only as a drone control system but also as a foundational platform for broader robotics and autonomous systems. By providing a structured and maintainable base layer, it enables developers and manufacturers to focus on application-driven innovation without being constrained by low-level system complexities.

We envision a future driven by autonomous systems capable of human-independent decision-making. Achieving this vision requires control over the foundational layers of the technology stack. This work represents an effort to move beyond assembly-based development toward building complete systems with full ownership of architecture and behavior.

The remainder of this report is structured as follows. Chapter~\ref{chap:preliminary} discusses existing systems and their limitations. Chapter~\ref{chap:architecture} presents the proposed system architecture. Chapter~\ref{chap:hardware} describes the hardware platform, followed by Chapters~\ref{chap:navhal} and \ref{chap:vaios}, which detail the design of NavHAL and VAIOS. Finally, Chapter~\ref{chap:vayu} presents the Vayu flight control stack.